\begin{example}[Refuted: bounding rational chains by local degree and different]
\label{D:fixed-different-unbounded-rational-chain}
Let \(k\) be algebraically closed of characteristic two. For \(m\geq1\),
put \(N=2m+1\) and
\[
 X_m=\operatorname{Spec}
 k[x,y,z]/(z^2+xz+y^N),\qquad
 \pi_m:X_m\longrightarrow\mathbb A^2_{x,y}.
\]
These are finite flat, generically separable degree-two normal
coverings, etale away from the fixed divisor \(x=0\).
At its generic point there is one extension of the valuation, with
\[
 e=1,\qquad f_{\mathrm{insep}}=2,\qquad
 f_{\mathrm{sep}}=1,\qquad \delta=1.
 \tag{R1}
\]
Nevertheless the minimal resolution at the point over the origin
has a chain of \(2m\) smooth rational curves, each of square \(-2\);
the absolute determinant of their intersection matrix is \(2m+1\).
The reduced curve over \(x=0\) has local normalization defect
\(\delta_{\mathrm{curve}}=m\).
Thus a fixed smooth base germ, fixed covering degree, and fixed
divisorial different do not bound the length or determinant of
these rational chains. This statement is local along a nonproper
branch curve; it does not supply a covering in the contracted-point
setup of the finite-generation problem.
\end{example}

\begin{proof}
The monic equation in \(z\) makes the covering free of rank two.
Its derivative in \(z\) is \(x\), so the map is etale on \(D(x)\)
and is generically separable.
The change of variables \(u=x+z,\ v=z\) gives
\[
 k[X_m]=k[u,v,y]/(uv+y^N).
 \tag{R2}
\]
This is a domain, for example by viewing the primitive polynomial
\(uv+y^N\) as a linear polynomial in \(u\) over the unique
factorization domain \(k[v,y]\).
Its Jacobian has entries \(v,u,y^{N-1}\), so its only singular
point is the origin. A hypersurface is Cohen--Macaulay, and
Serre's criterion proves normality
\cite[Tag 031S]{StacksRationalChainB}.

The quotient by \(x\) is
\(k[y,z]/(z^2+y^N)\), a domain with normalization
\[
 y=s^2,\qquad z=s^N.
\]
At the generic point its residue field is \(k(s)\), purely
inseparable of degree two over \(k(y)\). The height-one prime
\((x)\) in the normal ring is principal, so \(x\) is a uniformizer
there and \(e=1\). The monogenic different is generated by
the derivative \(x\), by
\cite[Example 49.12.5, tag 0BWB]{StacksRationalChainB}.
Its valuation is one. This proves (R1).
For the curve at the origin, completion gives
\(k[[s^2,s^N]]\subset k[[s]]\). Its missing exponents are exactly
\(1,3,\ldots,2m-1\), so the quotient has \(k\)-dimension \(m\).

We give the resolution explicitly, avoiding any restriction on
the characteristic in a classification theorem.
For \(i=0,\ldots,N-1\), take a smooth chart with coordinates
\(a_i,b_i\) and maps
\[
 u=a_i^{\,i+1}b_i^{\,i},\qquad
 v=a_i^{\,N-i-1}b_i^{\,N-i},\qquad y=a_i b_i.
 \tag{R3}
\]
Neighbouring charts are glued, where \(b_i\ne0\), by
\[
 a_{i+1}=b_i^{-1},\qquad b_{i+1}=a_i b_i^2.
 \tag{R4}
\]
The three formulas (R3) agree on these overlaps.

This glued map is proper. To verify this directly, blow up the
ideal \((u,y)\) on (R2). On the \(u\)-chart put \(y=ua\);
the equation gives \(v=u^{N-1}a^N\), a smooth affine plane.
On the \(y\)-chart put \(u=yb\); the remaining equation is
\(bv=y^{N-1}\). Repeat the same construction on that chart.
The next ideal \((b,y)\) is the unit ideal on its overlap with
the already smooth chart, so it glues to the unit ideal there.
Induction gives exactly (R3)--(R4). Every step is a blowup of
a coherent ideal, hence is projective
\cite[Tag 02NS]{StacksRationalChainB}; it is an isomorphism
away from the original singular point. The resulting smooth
surface is therefore a proper resolution.

For \(j=1,\ldots,N-1\), the exceptional curve \(C_j\) is
\(a_{j-1}=0\) in chart \(j-1\) and \(b_j=0\) in chart \(j\).
The transition \(a_j=b_{j-1}^{-1}\) makes it a projective line.
The normal-coordinate transition in (R4) is
\(b_j=a_{j-1}b_{j-1}^2\); hence its normal bundle is
\(\mathcal O_{\mathbf P^1}(-2)\) and \(C_j^2=-2\).
In the intermediate charts the two adjacent exceptional curves
are the coordinate axes; they meet transversely once.
Nonadjacent curves do not meet. These are all exceptional
components, and none is a rational \((-1)\)-curve, so the
resolution is minimal.

The determinant \(D_t\) of the positive tridiagonal matrix with
diagonal \(2\) and adjacent entries \(-1\) satisfies
\(D_0=1,D_1=2,D_t=2D_{t-1}-D_{t-2}\).
Thus \(D_t=t+1\); taking \(t=N-1=2m\) proves the asserted
determinant. The same explicit formal charts give the resolution
after completing at the origin, with the same exceptional
curves and local invariants.
\end{proof}

\begin{remark}[What global contraction data must still control]
The normalized residual genus in this family is zero and its
separable residual different is zero. Its unbounded contribution
is the arithmetic defect of the nonnormal residual curve.
A global estimate on a proper exceptional fibre may constrain
that defect; the local calculation does not decide this.
It does show that a determinant estimate cannot follow from the
divisorial \(\delta/e\) bound and the degree alone.
\end{remark}

